\documentclass[11pt]{article}
\usepackage[margin=25mm]{geometry}
\usepackage{fontspec,unicode-math}
\setmainfont{Latin Modern Roman}
\setmathfont{Latin Modern Math}
\usepackage{amsmath,amssymb,mathtools}
\usepackage{xurl,hyperref}
\hypersetup{colorlinks=true,urlcolor=blue,linkcolor=blue}
\setcounter{secnumdepth}{0}
\setlength{\parindent}{0pt}
\setlength{\parskip}{5pt}
\setlength{\emergencystretch}{3em}
\newcommand{\sref}[2]{\hyperlink{source:#1:#2}{[S#2]}}
\newcommand{\cataloguescope}[1]{\paragraph{Source question.}#1}
\newcommand{\floor}[1]{\left\lfloor#1\right\rfloor}
\begin{document}
\section{Better-Than-Two Approximation for Multiplicative Connectivity Augmentation}\label{statement}

\subsection{Definitions and mathematical statement}

\paragraph{Shared notation.}Standard number systems are denoted by $\mathbb N,\mathbb Z,\mathbb Q,\mathbb R,\mathbb C$; index ranges, computational models and specialized notation follow the cited problem source and the definitions below.


In $\alpha$-multiplicative connectivity augmentation the input is a $k$-edge-connected graph and a set of candidate links with nonnegative costs. One selects minimum-cost links so that the augmented graph is at least $\lceil\alpha k\rceil$-edge-connected. Edge connectivity is the minimum number of edges in a nontrivial cut. An approximation factor below two means a polynomial-time algorithm with cost at most $\rho$ times optimum for some $\rho<2$.
\paragraph{Statement recorded in the supplied source.}
For every fixed $\alpha>1$, does $\alpha$-multiplicative connectivity augmentation admit an approximation factor strictly below $2$? \sref{statement}{2}
\paragraph{Scope qualification.}
The improvement may depend on fixed $\alpha$ but must be uniform over the input instances. Feasibility assumes enough available links to achieve the target.
\paragraph{Source status.}
The supplied catalogue labels this record partially solved. Its source-status metadata was checked against the supplied snapshot on 2026-10-11. This import does not independently establish current open status or mathematical correctness. The catalogue literature assessment is dated 2026-08-22.
\subsection{Short English statement}
Improve the factor-two approximation barrier for increasing edge connectivity by a fixed multiplicative factor.
\subsection{Sources}
\begin{itemize}
\item[\hypertarget{source:statement:1}{S1}] Ulam AI and the UnsolvedMath Contributors, UnsolvedMath: A Curated Collection of Open Mathematics Problems, catalogue entry. Catalogue curation and metadata: CC BY 4.0. The supplied source question is retained; definitions and the English summary are editorial additions.
\url{https://huggingface.co/datasets/ulamai/UnsolvedMath}.
\item[\hypertarget{source:statement:2}{S2}] Combinatorial Optimization (2022), Oberwolfach Reports, DOI 10.4171/owr/2021/53, pp. 2893--2954
\url{https://doi.org/10.4171/owr/2021/53}.
\item[\hypertarget{source:statement:3}{S3}] Combinatorial Optimization, Oberwolfach Reports 53/2021, primary report PDF
\url{https://ems.press/content/serial-article-files/46931}.
\textit{Underlying Oberwolfach Reports material retains its source-specific licence, including CC BY-SA 4.0 where applicable; see THIRD\_PARTY\_NOTICES.md.}
\end{itemize}
\paragraph{Preparation and review.}
Original question attributed to the cited source and UnsolvedMath contributors. Codex (GPT-6; exact serving revision not exposed) assisted with definitions, English summaries, source inspection and formatting. No new human mathematical authorship or independent proof verification is asserted. Catalogue identity checked against the complete supplied registry and existing repository statements; known flagged repetitions and obvious identical conjectures excluded. Full mathematical-equivalence/literature review has not been performed. Statement preparation and final review followed contribution prompts C and D; automated checks verify structure and identity, not open status.
\label{end:statement}
\end{document}
